% 表格风格对比样张（4 种方案）
\documentclass[11pt,a4paper]{ctexart}
\usepackage[margin=2.4cm]{geometry}
\usepackage{tabularray}
\usepackage{xcolor}
\usepackage{booktabs}
\usepackage{float}
\definecolor{ink}{HTML}{1F3864}    % 墨蓝
\definecolor{hdrblue}{HTML}{DCE6F1}% 淡蓝灰表头
\definecolor{rowalt}{HTML}{F5F7FA} % 极淡隔行
\definecolor{linegray2}{HTML}{C9CFD6} % 更浅细线

\begin{document}

\section*{方案 A —— 现行 booktabs 三线表（基线）}
\begin{table}[H]\centering
\caption{方案 A：三线表（基线）}
\begin{tabular}{@{}p{6.2cm}p{2.6cm}p{2.4cm}p{3.2cm}@{}}
\toprule
原始内容或位置 & 提取结果 & 单位或上下文 & 当前状态 \\
\midrule
1999ApJ…523..328N 第 13 页正文：“Hipparcos 视差误差相关性”上下文 & 130.7 $\pm$ 11.1 & pc & 正确 \\
2005A\&A…429.645S 第 1 页表格（距离与金属丰度测定表） & $-$0.034 $\pm$ 0.024 & dex & 正确 \\
2018MNRAS…473.4612K 表格单元格 “7.48 $\pm$ 0.03” & 7.48 & mas & 已统一（清理不确定度） \\
MWSC 星团目录 catalog 行 0305 原始列 logt & 8.15 & log(yr)（原始写法保留） & 待核对（与论文 Myr 写法不同） \\
2006A\&A…446..611F 第 2 页 “星团参数汇总表” & 75-100 & Myr（区间保留原样） & 待核对（区间按设计保留） \\
\bottomrule
\end{tabular}
\end{table}

\clearpage
\section*{方案 B —— 表头浅底 + 细线 + 隔行条纹（无竖线，最克制）}
\begin{table}[H]\centering
\caption{方案 B：浅灰表头 + 细横线 + 隔行条纹}
\begin{tblr}{colspec={X[3.2]X[1.1]X[1.2]X[1.6]}, row{1}={bg=hdrblue,font=\bfseries}, row{odd[2-Z]}={bg=rowalt}, hline{1,Z}={0.8pt,ink}, hline{2}={0.4pt,ink}, hline{3-Y}={0.2pt,linegray2}}
原始内容或位置 & 提取结果 & 单位或上下文 & 当前状态 \\
1999ApJ…523..328N 第 13 页正文：“Hipparcos 视差误差相关性”上下文 & 130.7 $\pm$ 11.1 & pc & 正确 \\
2005A\&A…429.645S 第 1 页表格（距离与金属丰度测定表） & $-$0.034 $\pm$ 0.024 & dex & 正确 \\
2018MNRAS…473.4612K 表格单元格 “7.48 $\pm$ 0.03” & 7.48 & mas & 已统一（清理不确定度） \\
MWSC 星团目录 catalog 行 0305 原始列 logt & 8.15 & log(yr)（原始写法保留） & 待核对（与论文 Myr 写法不同） \\
2006A\&A…446..611F 第 2 页 “星团参数汇总表” & 75-100 & Myr（区间保留原样） & 待核对（区间按设计保留） \\
\end{tblr}
\end{table}

\clearpage
\section*{方案 C —— 细灰全框 + 淡蓝表头（边框齐全但轻）}
\begin{table}[H]\centering
\caption{方案 C：细灰全框 + 淡蓝表头}
\begin{tblr}{colspec={X[3.2]X[1.1]X[1.2]X[1.6]}, row{1}={bg=hdrblue,font=\bfseries}, row{odd[2-Z]}={bg=rowalt}, hline{1,Z}={0.9pt,ink}, hline{2}={0.4pt,ink}, hline{3-Y}={0.2pt,linegray2}, vline{1,Z}={0.9pt,ink}, vline{2-Y}={0.2pt,linegray2}}
原始内容或位置 & 提取结果 & 单位或上下文 & 当前状态 \\
1999ApJ…523..328N 第 13 页正文：“Hipparcos 视差误差相关性”上下文 & 130.7 $\pm$ 11.1 & pc & 正确 \\
2005A\&A…429.645S 第 1 页表格（距离与金属丰度测定表） & $-$0.034 $\pm$ 0.024 & dex & 正确 \\
2018MNRAS…473.4612K 表格单元格 “7.48 $\pm$ 0.03” & 7.48 & mas & 已统一（清理不确定度） \\
MWSC 星团目录 catalog 行 0305 原始列 logt & 8.15 & log(yr)（原始写法保留） & 待核对（与论文 Myr 写法不同） \\
2006A\&A…446..611F 第 2 页 “星团参数汇总表” & 75-100 & Myr（区间保留原样） & 待核对（区间按设计保留） \\
\end{tblr}
\end{table}

\clearpage
\section*{方案 D —— 墨蓝表头白字 + 细灰内线（数据仪表风）}
\begin{table}[H]\centering
\caption{方案 D：墨蓝表头 + 白字 + 细灰内线}
\begin{tblr}{colspec={X[3.2]X[1.1]X[1.2]X[1.6]}, row{1}={bg=ink,fg=white,font=\bfseries}, row{odd[2-Z]}={bg=rowalt}, hline{1,Z}={0.9pt,ink}, hline{2}={0.5pt,ink}, hline{3-Y}={0.15pt,linegray2}, vline{1,Z}={0.9pt,ink}, vline{2-Y}={0.15pt,linegray2}}
原始内容或位置 & 提取结果 & 单位或上下文 & 当前状态 \\
1999ApJ…523..328N 第 13 页正文：“Hipparcos 视差误差相关性”上下文 & 130.7 $\pm$ 11.1 & pc & 正确 \\
2005A\&A…429.645S 第 1 页表格（距离与金属丰度测定表） & $-$0.034 $\pm$ 0.024 & dex & 正确 \\
2018MNRAS…473.4612K 表格单元格 “7.48 $\pm$ 0.03” & 7.48 & mas & 已统一（清理不确定度） \\
MWSC 星团目录 catalog 行 0305 原始列 logt & 8.15 & log(yr)（原始写法保留） & 待核对（与论文 Myr 写法不同） \\
2006A\&A…446..611F 第 2 页 “星团参数汇总表” & 75-100 & Myr（区间保留原样） & 待核对（区间按设计保留） \\
\end{tblr}
\end{table}

\end{document}
